\subsection{Применение интеллектуальных инструментов в инженерном проектировании}
\label{sec:development:ai}

Под интеллектуальными инструментами здесь понимаются большие языковые модели, которые принимают постановку задачи на естественном языке и возвращают текст, код или описание схемы. При разработке системы применялась модель \textit{Claude} в среде \textit{Claude Code}, работающая с файлами проекта. Инструмент использовался как помощник на трёх этапах: при проектировании модели данных, при написании кода и при подготовке графического материала. Решение о применении каждого предложенного варианта принимал разработчик после проверки.

\subsubsection{Порядок работы с инструментом}

Работа строилась по одной схеме. Задача формулировалась вместе с контекстом: фрагментом схемы базы данных, текстом запроса или сообщением об ошибке. Полученный ответ не принимался как готовый результат, а проверялся: код~-- запуском тестов, запросы~-- планом выполнения, схема данных~-- сверкой с требованиями подраздела~\ref{sec:analysis:specification}. Ответы, которые не удавалось воспроизвести или обосновать, отбрасывались.

Такой порядок выбран из-за особенности инструмента: модель порождает синтаксически верный и правдоподобный результат, но не проверяет его на данных системы. Ошибка в ответе выглядит так же убедительно, как верное решение, поэтому проверка обязательна для каждого предложения.

\subsubsection{Решённые задачи}

Задачи, при решении которых применялся инструмент, сведены в таблицу~\ref{table:development:ai}.

\begin{table}[ht]
\caption{Задачи, решённые с применением интеллектуальных инструментов}
\label{table:development:ai}
\centering
	\begin{tabular}{|>{\raggedright}m{0.26\textwidth}
	                |>{\raggedright}m{0.31\textwidth}
	                |>{\raggedright\arraybackslash}m{0.31\textwidth}|}
	\hline
	\centering Задача & \centering Результат работы инструмента & \centering\arraybackslash Что потребовало проверки \\
	\hline Проектирование модели данных & Черновой состав сущностей и связей по словесному описанию предметной области & Мощность связей, состав ключей, выделение совпадения в отдельную сущность \\
	\hline Подбор индексов & Перечень индексов под запросы цикла опроса & Соответствие плану выполнения запроса, поиск избыточных индексов \\
	\hline Разбор ответа площадки & Черновой код разбора полей объявления и характеристик категории & Поведение на объявлениях без цены и без описания \\
	\hline Миграции схемы & Заготовки миграций с прямым и обратным переходом & Правила удаления связанных записей, порядок применения \\
	\hline Тесты & Наборы проверок для отбора объявлений и лимитов уведомлений & Полнота граничных случаев \\
	\hline Графический материал & Исходные описания схем и диаграмм для \textit{Graphviz} и \textit{PlantUML} & Соответствие ГОСТ~19.701–90 и фактическому коду \\
	\hline
	\end{tabular}
\end{table}

\subsubsection{Ошибки, найденные при проверке}

Проверка результатов выявила два дефекта, которые относятся непосредственно к базе данных и были исправлены отдельными миграциями.

1. Правила удаления связанных записей существовали только в коде приложения, а в схеме базы данных отсутствовали. Удаление товара, по которому уже отправлено уведомление, завершалось нарушением ссылочной целостности. Миграция добавила правила в саму схему: для товаров, совпадений и значений справочника задано каскадное удаление, а ссылка журнала на совпадение обнуляется, поскольку запись журнала хранит собственную копию цены и идентификатора объявления и должна пережить удаление совпадения.

2. Среди индексов таблицы объявлений оказался избыточный: отдельный индекс по площадке дублировал составной индекс по паре «площадка~-- внешний идентификатор», который начинается с того же столбца и обслуживает те же запросы. Избыточный индекс удалён, потому что он замедлял вставку объявлений и занимал место, не ускоряя ни один запрос.

Оба дефекта не были обнаружены ни тестами, ни разбором кода: первый проявляется только при удалении товара с историей уведомлений, второй вообще не влияет на правильность результата. Нашла их сверка схемы с перечнем запросов цикла, выполненная вручную.

\subsubsection{Оценка применимости}

Инструмент даёт выигрыш на задачах, где ответ проверяется быстро и однозначно: черновик кода разбора ответа площадки, заготовка миграции, набор тестов, описание диаграммы для \textit{Graphviz}. Время на такие задачи сокращается заметно, потому что проверка занимает меньше времени, чем написание с нуля.

На задачах, где решение зависит от особенностей площадки и характера нагрузки, выигрыш оказался небольшим. Ограничение частоты запросов, стратегия пауз после отказа площадки, порог выгодности и правила повторной отправки уведомления подбирались по наблюдениям за работой системы, а не по предложенному варианту. Модель не располагает сведениями о поведении конкретной площадки и о распределении цен в категориях, поэтому её предложения в этой части оставались общими.

Отсюда следует правило, которого придерживались при разработке: интеллектуальный инструмент применим там, где результат проверяем формально, и не применим там, где решение опирается на данные, которых у инструмента нет. Ответственность за принятое решение остаётся на разработчике в обоих случаях.
